\section{Discussion}
\label{sec:discussion}

\subsection{Why there was no crossover}

The idea of a crossover rests on an asymmetry: a monolith can only be scaled as a whole, so at some load giving
resources to just the overloaded part should win. Under a fixed budget the asymmetry runs the other way, because
decomposition cannot add resources; it can only divide them. The treatments separate what dividing them cost.

Part of the gap is fragmentation. Under the equal split the edge service exhausted its share while the rest of the
budget stood idle: at 64 clients C used \xo{budgetuse/browse/C}\% of its six cores and the monolith
\xo{budgetuse/browse/A}\%. The edge service's CFS throttling also made tail latency erratic at average utilisations
of \xo{hot/browse/C/500/util}--\xo{hot/browse/B/2200/util}\%. Allocating the budget by demand removed most of this and
raised the decomposed configurations' capacity relative to the monolith by a factor of
\xo{alloc/gainmin}--\xo{alloc/gainmax}.

The rest is the work that distribution adds to every request, \xo{perhop/browse/cpu/min}--\xo{perhop/browse/cpu/max}\,ms
of CPU per remote call without a mesh. No division of a fixed budget recovers it, because even a perfect allocation
is bounded by the ratio of the CPU each request needs. How much that matters depends on how much work a request
carries. With microsecond business logic the added work exceeded the logic itself. With 1\,ms in every module
operation it became a fraction of it, and the proportional allocation brought B to \xo{workprop/peakpct/max}\% of
the monolith's capacity. The cloud environment shows the same structure with different weights. Its monolith was
faster at low load, so the same few tenths of a millisecond per call weighed more in the latency ratio, while its
database, reached across the network, saturated earlier and left a smaller capacity gap.

None of this means that decomposition cannot win on performance. It means that its case depends on resources the
monolith cannot use, such as a larger budget or one that follows the load. Berry et al.\ found the outcome to depend
on the scaling strategy~\cite{Berry2024}, and Blinowski et al.\ found that scaling out beyond some number of
instances degraded performance~\cite{Blinowski2022}.
Autoscaling under a shared ceiling did not supply that case either. The decomposed configurations' autoscalers
added replicas to the services that needed them, as the argument predicts, but each replica brought its own cold
JVM and its share of the per-call cost, and the monolith served more load on one replica than they served at the
ceiling.

\subsection{Implications for practitioners}

\paragraph{Price each synchronous call.} A remote call cost a few tenths of a millisecond of transit, more than its
serialisation, and a few tenths of a millisecond of CPU. A sidecar mesh doubled the transit locally and tripled it
in the cloud. The price can be measured cheaply, with a \texttt{Server-Timing} header and a load generator that
records it, and it should be set against the work a call carries before deciding how finely to cut.

\paragraph{Size the edge service, not the service count.} The capacity of every decomposed configuration was its
edge service's CPU share. That service's throttled-period counter warns earlier than its utilisation, and sizing
services to their measured demand raised capacity relative to the monolith \xo{alloc/gainmin}--\xo{alloc/gainmax}
times without adding CPU.

\paragraph{Prefer coarse splits, and count the mesh's CPU.} When organisational reasons drive a split, the coarsest
split that serves them costs least. B kept more capacity than C under every allocation and in both environments, up
to three times as much, which puts numbers on the advice against services that are too small~\cite{Taibi2018}. On a
sidecar mesh the proxies used more CPU than the application. Sidecar-less data planes are reported to cost much
less~\cite{BremlerBarr2025} and should be measured in the same way.

\subsection{Lessons for testbed builders}

Building the protocol changed it. The original latency decomposition subtracted span timestamps taken on different
hosts and therefore depended on clock synchronisation; the in-band accounting does not. The original planned to fit
the USL to open-loop data, a model of a closed system applied to an open one. Its mesh configuration also added
asynchronous messaging, which would have changed two variables at once. Beyond the protocol, several defaults of
the software stack bias a comparison of this kind, and others building such testbeds will meet them:
\begin{itemize}
\item The JVM chooses its garbage collector from the CPU limit, so under an equal split the fine-grained services
got the serial collector and the monolith G1. The collector must be pinned.
\item A Spring Boot JVM needs about 300\,MiB at idle. An equal memory split leaves small services open to OOM kills
under overload, so memory needs a floor and is better reported as a cost than equalised.
\item Unbounded concurrency (virtual threads, unlimited connections) turns overload into memory growth and OOM
kills. Bounding it makes overload queue in the kernel instead.
\item The JDK HTTP client does not retry a POST that races with a server-side connection close, so servers must not
close keep-alive connections after a request limit during measurement.
\item Against a remote cluster, \texttt{kubectl exec} streams can end before all output has arrived without
reporting an error. Every transferred result file must be checked.
\end{itemize}

\subsection{Future work}

The campaign runner also targets bare-metal clusters. Runs on x86 bare metal would separate virtualisation and
cloud-network effects from those of the architectures. The autoscaling treatment can be extended with warmed
replicas, other utilisation targets and a database tier that scales as well. Further variants are gRPC instead
of JSON over HTTP/1.1, a sidecar-less mesh, and a workload that mixes operations.

\section{Conclusion}
\label{sec:conclusion}

We compared a monolith with three decompositions of identical code, on a local cluster and on Amazon EKS, under an
equal CPU budget and a demand-proportional allocation and at two levels of service demand.
We also added autoscalers under a shared CPU ceiling in the cloud.
Decomposition added latency at low load: the p99 of the decomposed configurations was \xo{h1/min}--\xo{h1/max}
times the monolith's locally and \xo{eks/h1/min}--\xo{eks/h1/max} times in the cloud, mostly inter-service transit.
It cut capacity to \xo{peakpct/min}--\xo{peakpct/max}\% (local) and \xo{eks/peakpct/min}--\xo{eks/peakpct/max}\%
(cloud) of the monolith's, because the edge service exhausted its CPU share while the rest of the budget stayed
idle. Allocating the budget by demand raised this to at most \xo{prop/peakpct/max}\%, and with 1\,ms of CPU in every
module operation to at most \xo{workprop/peakpct/max}\%. A sidecar mesh added
\xo{mesh/perhop/min}--\xo{eks/mesh/perhop/max}\,ms per call and more CPU than the application. With autoscalers the
monolith sustained \xo{hpa/sustained/browse/A} browse requests per second on one replica, the decomposed
configurations at most \xo{hpa/sustained/browse/B} at the ceiling. No decomposed configuration overtook the monolith
in any setting.

For a team deciding how finely to split, these measurements put a price on each synchronous call and show where it
is paid. They do not show a load at which decomposition pays for itself in performance; its case rests on the other
benefits it brings, and the price is smallest when services are coarse and do substantial work. The testbed, data
and analysis are published so that the same measurements can be repeated on other infrastructure.
